\chapter{术语、符号与前置知识}

\section{本章作用}

本章用于建立后续章节共享的术语框架。Pebble 的讨论对象横跨产品体验、关系数据、交互反馈与分析服务，如果读者在进入架构与方法章节前没有先区分这些对象，就容易把系统的设计目标、已观察能力与用户感知层面的体验描述混写在一起。因此，本章先给出后文最常出现的概念，并说明它们在本文中的含义边界。

对于第一次接触 Pebble 的读者，建议把本章理解为一份最小词汇表。后文若出现相同术语，默认沿用本章定义；若某一章节对术语做了更细的展开，应视为在本章定义基础上的局部细化，而不是重新定义对象。

\section{核心术语总表}

表~\ref{tab:pebble-terms} 汇总了后文最重要的术语。它的目的在于确保读者在阅读产品主张、系统设计和验证部分时，对同一个词保持同一种理解，而不在这一章提前展开全部实现细节。

\begin{table}[H]
  \centering
  \caption{Pebble 核心术语与本文中的使用边界。}
  \label{tab:pebble-terms}
  \begin{tabularx}{\textwidth}{>{\raggedright\arraybackslash}p{2.8cm} >{\raggedright\arraybackslash}X >{\raggedright\arraybackslash}X}
    \toprule
    术语 & 本文中的含义 & 阅读时应注意的边界 \\
    \midrule
    关系上下文 & 与某个关系对象相关的背景信息、互动历史、预期目标与场景补充 & 它是 Pebble 解释对话与生成建议的中心条件，不应被理解为附属备注 \\
    对话解码 & 对当前输入话语做表层语义、潜台词与回应候选的分析过程 & 它解决“先理解”问题，不单独承担长期训练与状态恢复职责 \\
    灰岩回复 & 尽量降低情绪刺激、减少解释与对抗升级的回应策略 & 它是策略风格，不等于冷漠或机械回避一切沟通 \\
    模拟陪练场 & 通过场景化对话练习回应方式，并获得即时反馈的交互环境 & 它强调训练与复盘，不直接替代真实关系中的最终决策 \\
    急救呼吸 & 在情绪负荷过高时提供短时稳态支持的生理调节模块 & 它处理的是状态恢复，不负责完成语言分析本身 \\
    关系档案 & 面向具体关系对象持续维护的结构化信息集合 & 它既服务即时分析，也为长期洞察和策略适配提供基础 \\
    已观察事实 & 已经在当前工作区、页面或接口中得到直接观察的信息 & 该层只描述真实存在的能力，不自动外推为未来保证 \\
    design intent & 报告中关于目标结构、产品方向与系统设计的规划性描述 & 该层表达的是要建成什么，不等于当前已经全部实现 \\
    \bottomrule
  \end{tabularx}
\end{table}

\section{需要先建立的直觉}

在 Pebble 的语境中，读者至少需要建立三个直觉判断。

\begin{itemize}
  \item 语言分析不是孤立任务。同一句话在不同关系中可能意味着完全不同的压力结构，因此 Pebble 不能只围绕文本表面做响应。
  \item 高质量支持不等于一次回答。用户往往需要先理解，再练习，再恢复状态，并在后续互动中继续积累上下文。
  \item 系统价值来自协同，而非单模块堆叠。解码、陪练、稳态恢复与关系档案只有在同一条支持链路中协同工作时，Pebble 的整体主张才成立。
\end{itemize}

这三个直觉决定了后文的阅读方式。若读者始终把 Pebble 理解成“一个分析文本的页面”，那么第五章之后的很多设计选择会显得过度复杂；若读者接受 Pebble 是一个围绕复杂关系组织支持链路的系统，那么这些设计选择就会表现为必要的结构条件。

\section{读者需要掌握的最小前置知识}

理解本报告并不要求读者具备心理学、临床医学或复杂机器学习背景，但仍需掌握若干最小前提。

\begin{itemize}
  \item 读者应理解“上下文会改变同一句话的解释结果”这一基本事实。
  \item 读者应接受产品设计中的边界意识，即系统可以提供支持，但不替代用户做价值判断。
  \item 读者应能区分界面层、服务层与数据层这三类对象，避免把页面文案、分析逻辑和关系数据视为同一种设计元素。
  \item 读者应能区分规划性描述与已验证能力，这一点对理解本报告中的 design intent 与已观察事实尤为重要。
\end{itemize}

如果读者具备以上前置知识，后续章节中的系统边界、交互流程和验证策略就可以被直接理解；若缺少这些前提，则建议在继续阅读前先回看第一章与第二章，重新确认 Pebble 的问题域、目标边界与核心主张。
